% Source: sources/returns/S6_04_MAPS_MANIFOLDS/sections/manifolds.tex:4291–4396
\paragraph{Theorem D.2.21 (Proved).}
Use the uncorrected degree-one data $P=G,Q=dG$ and the simple zero at the
marked $I_1$ in the following five rows. Keep the indicated original
additive filling and replace the other additive fibre by its finite-good
logarithmic model of order $m=qk$, with invariant power vector
$n\ell+r\delta$, where $k>0$, $r\in\mathbb Z$ and $\gcd(n,m)=1$.
\[
\begin{array}{c|c|c|r|r|r|r}
\mathrm{OS/profile}&d&\mathrm{finite\ log}\ (q)&\mathrm{original}&b&h&c\\\hline
49/III&6&IV^*\ (3)&III&1&3&2\\
49/III&6&III\ (4)&IV^*&6&2&3\\
50/IV&12&IV\ (3)&I_1^*&5&3&4\\
43/II&2&II\ (6)&III^*&7&1&2\\
46/II&4&II\ (6)&I_3^*&7&1&4
\end{array}
\]
Then $n\ne0$, $\pi_1(X)=H_1(X)=\mathbb Z/|n|$, and
\[
 H_2(X)=\mathbb Z^b,\qquad H_3(X)=\mathbb Z^{2b},\qquad
 H_4(X)=\mathbb Z^b\oplus\mathbb Z/|n|,\qquad H_5(X)=0.
\]
The endpoint groups are $\mathbb Z$, and $\chi(X)=2$.
Every member with $n=\pm1$ is oriented diffeomorphic to
\[
                 \#^b(S^2\times S^4)\ \#\ \#^b(S^3\times S^3).
\]
In particular none of these five simply connected families is $S^6$.

\emph{Proof.} The component data, original local models and geometric
finite-good logarithmic replacements are those already proved in
D.2.20 and D.2.5--D.2.12. We retain their markings rather than replacing
an original additive fibre by its invariant torus cohomology. Order the
three monodromies cyclically so that the finite-good one is first, and
write them $A_0,A_1,A_2$. Their product is the identity. The two remaining
full fixed lattices provide four primitive normal-torus vectors $B$.
If $\ell$ is the finite invariant vector of $\gamma$-coefficient one,
then $K=[B,j_0\ell]$ is a primitive integral basis of the rank-five
kernel of the graph differential. The exact integer inverses, monodromies,
component projections and complete matrices are saved in
\texttt{code/manifolds\_mixed\_additive.py} and its JSON certificate.

The regular degree-two coinvariants are freely generated by $s$, detected
by $\omega=e_1^*e_2^*-d\delta^*\gamma^*$. The original additive cap has
$r_0=b+1$ components and the actual principal $T^2$-bundle model of
D.2.20. Its local free defect is $r_0-2=b-1$. Its full fibre kernel maps
to $h\mathbb Zs$, with $h$ in the table. The one-component toric cusp
has no free defect. The finite cap has no free defect either.

Let $\beta$ generate the finite vertical coinvariant lattice and write
$\delta=q_\delta\beta$. Then
\[
 N_m\beta=(m/q_\delta)\delta,\qquad N_m\ell=m\ell,
 \qquad -\omega(\beta,\ell)=c=d/q_\delta,\quad
 -\omega(\ell,\delta)=-d.
\]
Choose $m\alpha+n\zeta=1$. The actual finite normal-circle tori are
\[
 \binom{-\omega(\beta,n\ell+r\delta)}{Lj_0(N_m\beta)},\qquad
 \binom{-\alpha\omega(\ell,n\ell+r\delta)}
 {Lj_0(\alpha N_m\ell+\zeta(n\ell+r\delta))}.
\]
After eliminating the four primitive columns $B$, the entire remaining
integral cokernel has the presentation
\[
                 \begin{pmatrix}h&cn&-dr\alpha\\0&0&1\end{pmatrix}.
\]
In each row $\gcd(h,cn)=1$: the nontrivial possibilities use respectively
$3\nmid n$, $2\nmid n$, $3\nmid n$; the last two have $h=1$.
Thus this cokernel vanishes for every allowed integer input, without an
order bound. The companion code verifies every constant primitive lattice
and $1,068$ full $6\times7$ Smith matrices as an independent certificate.

For completeness, the first-degree Mayer--Vietoris source has rank nine
and target rank eight: each boundary has rank three; the exterior has
rank three, while the finite, original and cusp caps have ranks two,
one and two. The original cap kills the elliptic and $\delta$ loops,
leaving $\gamma$. All original and cusp meridians vanish. The finite
power relation is therefore $n\gamma=0$. This proves
$\pi_1(X)=\mathbb Z/|n|$ directly and makes the rational first-degree
Mayer--Vietoris map onto. Its kernel is consequently free of rank one.
Together with the local defect $b-1$ and the zero cokernel above, this
proves $H_2(X)=\mathbb Z^b$. Integral duality, the universal coefficient
theorem and $\chi=2$ give the remaining groups as displayed.

When $n=\pm1$, the reduced finite fibre divisor is zero in integral
cohomology: the complement has fundamental group $\mathbb Z\gamma$,
and its primitive normal meridian maps to $\pm n\gamma$, so the divisor
residue sequence gives divisor order $|n|$. The canonical divisor differs
from a base pullback by a multiple of this reduced fibre; the full base
fibre is $m$ times the same divisor. Thus $c_1(X)=0$, including on the free
Mayer--Vietoris bridge class. The $\delta$ action is free on the original
principal bundle and on the finite quotient, since $\gcd(n,m)=1$.
Its only fixed component is the standard cusp sphere with normal weights
$+1,-1$. As in D.2.20 this sphere lies in the image of the regular fibre,
hence is a multiple of $s$ and is nullhomologous by the unimodular
presentation. The equivariant surgery, spin lift and Smale--Barden
argument of D.2.14 now give a free circle action on
$X\#(S^3\times S^3)$ with orbit
$\#^{b+1}(S^2\times S^3)$ and primitive Euler class.
Repeated Galaz-Garc\'ia--Reiser Theorem A and Theorem B(2), followed by
full Zhubr cancellation exactly as in D.2.16, give the asserted oriented
diffeomorphism type.\hfill$\square$


% Independently audited by root; integrated as D.2.22.
\paragraph{Theorem D.2.22 (Proved: finite initial data for sphere outputs).}
